% =========================================================
% Slides Beamer Metropolis 16:9 — Chapter 3 (Stationarity)
% Time Series Analysis — Aymen Ben Brik (aymen.benbrik@esprit.tn)
%
% Migrated from _archives/source-slides/Stationnarité.tex (1950 L,
% Madrid). Only stationarity + ACF/PACF + WN tests; AR/MA/ARMA goes
% to Ch.5, ADF/KPSS tests to Ch.4.
% =========================================================
\documentclass[aspectratio=169]{beamer}

% =========================================================
% Beamer preamble — Time Series Analysis module
% Theme : Metropolis, aspect ratio 16:9
% Author : Aymen Ben Brik (aymen.benbrik@esprit.tn)
% Esprit School of Business
%
% Include from each chapterX/slides.tex via :
%   % =========================================================
% Beamer preamble — Time Series Analysis module
% Theme : Metropolis, aspect ratio 16:9
% Author : Aymen Ben Brik (aymen.benbrik@esprit.tn)
% Esprit School of Business
%
% Include from each chapterX/slides.tex via :
%   % =========================================================
% Beamer preamble — Time Series Analysis module
% Theme : Metropolis, aspect ratio 16:9
% Author : Aymen Ben Brik (aymen.benbrik@esprit.tn)
% Esprit School of Business
%
% Include from each chapterX/slides.tex via :
%   \input{../preamble/slides-preamble.tex}
% AFTER the line \documentclass[aspectratio=169]{beamer}.
% =========================================================

\usepackage[utf8]{inputenc}
\usepackage[T1]{fontenc}
\usepackage[english]{babel}
\usepackage{lmodern}
\usepackage{amsmath, amssymb, mathtools, bm}
\usepackage{graphicx}
\usepackage{booktabs}
\usepackage{array}
\usepackage{multirow}
\usepackage{colortbl}
\usepackage{xcolor}
\usepackage{tikz}
\usetikzlibrary{positioning, arrows.meta, calc, shapes, fit, backgrounds, matrix, decorations.pathreplacing}
\usepackage{listings}
\usepackage[most]{tcolorbox}

% --- Theme ---
\usetheme{metropolis}
\setbeamertemplate{navigation symbols}{}
\metroset{block=fill}

% --- Esprit colour palette ---
\definecolor{espritBlue}{RGB}{30, 60, 110}
\definecolor{espritAccent}{RGB}{200, 80, 30}
\setbeamercolor{progress bar}{fg=espritAccent}
\setbeamercolor{title separator}{fg=espritAccent}
\setbeamercolor{frametitle}{bg=espritBlue}

% --- Code listings (Python + R, with UTF-8 literate) ---
\definecolor{codebg}{RGB}{248,248,248}
\definecolor{codeframe}{RGB}{220,220,220}
\definecolor{keyword}{RGB}{0,82,155}
\definecolor{comment}{RGB}{88,110,117}
\definecolor{string}{RGB}{133,153,0}

\lstdefinestyle{pythonstyle}{
  language=Python,
  basicstyle=\ttfamily\scriptsize,
  keywordstyle=\color{keyword}\bfseries,
  commentstyle=\color{comment}\itshape,
  stringstyle=\color{string},
  backgroundcolor=\color{codebg},
  frame=single,
  rulecolor=\color{codeframe},
  frameround=tttt,
  breaklines=true,
  showstringspaces=false,
  tabsize=2,
  literate=
    {é}{{\'e}}1 {è}{{\`e}}1 {ê}{{\^e}}1 {à}{{\`a}}1 {â}{{\^a}}1
    {ç}{{\c{c}}}1 {²}{{$^2$}}1 {°}{{${}^\circ$}}1
    {α}{{$\alpha$}}1 {β}{{$\beta$}}1 {γ}{{$\gamma$}}1
    {δ}{{$\delta$}}1 {θ}{{$\theta$}}1 {λ}{{$\lambda$}}1
    {μ}{{$\mu$}}1 {σ}{{$\sigma$}}1 {ε}{{$\varepsilon$}}1
    {ρ}{{$\rho$}}1 {τ}{{$\tau$}}1 {φ}{{$\phi$}}1
    {≤}{{$\leq$}}1 {≥}{{$\geq$}}1 {≈}{{$\approx$}}1
}
\lstdefinestyle{Rstyle}{
  language=R,
  basicstyle=\ttfamily\scriptsize,
  keywordstyle=\color{keyword}\bfseries,
  commentstyle=\color{comment}\itshape,
  stringstyle=\color{string},
  backgroundcolor=\color{codebg},
  frame=single,
  rulecolor=\color{codeframe},
  frameround=tttt,
  breaklines=true,
  showstringspaces=false,
  tabsize=2
}
\lstset{style=pythonstyle}

% --- Common macros ---
\input{../preamble/macros.tex}

% --- tcolorbox boxes for slides ---
\definecolor{boxTh}{RGB}{120, 30, 30}
\definecolor{boxProp}{RGB}{30, 100, 60}
\definecolor{boxDef}{RGB}{30, 60, 110}
\definecolor{boxEx}{RGB}{180, 120, 0}

\newtcolorbox{infobox}[1][]{enhanced, colback=espritBlue!5, colframe=espritBlue,
  fonttitle=\bfseries\small,
  title={\ifstrempty{#1}{}{#1}},
  boxrule=0.4pt, arc=2pt, left=4pt, right=4pt, top=2pt, bottom=2pt}
\newtcolorbox{warnbox}[1][]{enhanced, colback=espritAccent!5, colframe=espritAccent,
  fonttitle=\bfseries\small,
  title={\ifstrempty{#1}{}{#1}},
  boxrule=0.4pt, arc=2pt, left=4pt, right=4pt, top=2pt, bottom=2pt}
\newtcolorbox{thbox}[1][]{enhanced, colback=boxTh!5, colframe=boxTh,
  fonttitle=\bfseries\small,
  title={Theorem\ifstrempty{#1}{}{ -- #1}},
  boxrule=0.5pt, arc=2pt, left=4pt, right=4pt, top=2pt, bottom=2pt}
\newtcolorbox{propbox}[1][]{enhanced, colback=boxProp!5, colframe=boxProp,
  fonttitle=\bfseries\small,
  title={Proposition\ifstrempty{#1}{}{ -- #1}},
  boxrule=0.5pt, arc=2pt, left=4pt, right=4pt, top=2pt, bottom=2pt}
\newtcolorbox{defbox}[1][]{enhanced, colback=boxDef!5, colframe=boxDef,
  fonttitle=\bfseries\small,
  title={Definition\ifstrempty{#1}{}{ -- #1}},
  boxrule=0.5pt, arc=2pt, left=4pt, right=4pt, top=2pt, bottom=2pt}
\newtcolorbox{exbox}[1][]{enhanced, colback=boxEx!5, colframe=boxEx,
  fonttitle=\bfseries\small,
  title={Example\ifstrempty{#1}{}{ -- #1}},
  boxrule=0.5pt, arc=2pt, left=4pt, right=4pt, top=2pt, bottom=2pt}

% AFTER the line \documentclass[aspectratio=169]{beamer}.
% =========================================================

\usepackage[utf8]{inputenc}
\usepackage[T1]{fontenc}
\usepackage[english]{babel}
\usepackage{lmodern}
\usepackage{amsmath, amssymb, mathtools, bm}
\usepackage{graphicx}
\usepackage{booktabs}
\usepackage{array}
\usepackage{multirow}
\usepackage{colortbl}
\usepackage{xcolor}
\usepackage{tikz}
\usetikzlibrary{positioning, arrows.meta, calc, shapes, fit, backgrounds, matrix, decorations.pathreplacing}
\usepackage{listings}
\usepackage[most]{tcolorbox}

% --- Theme ---
\usetheme{metropolis}
\setbeamertemplate{navigation symbols}{}
\metroset{block=fill}

% --- Esprit colour palette ---
\definecolor{espritBlue}{RGB}{30, 60, 110}
\definecolor{espritAccent}{RGB}{200, 80, 30}
\setbeamercolor{progress bar}{fg=espritAccent}
\setbeamercolor{title separator}{fg=espritAccent}
\setbeamercolor{frametitle}{bg=espritBlue}

% --- Code listings (Python + R, with UTF-8 literate) ---
\definecolor{codebg}{RGB}{248,248,248}
\definecolor{codeframe}{RGB}{220,220,220}
\definecolor{keyword}{RGB}{0,82,155}
\definecolor{comment}{RGB}{88,110,117}
\definecolor{string}{RGB}{133,153,0}

\lstdefinestyle{pythonstyle}{
  language=Python,
  basicstyle=\ttfamily\scriptsize,
  keywordstyle=\color{keyword}\bfseries,
  commentstyle=\color{comment}\itshape,
  stringstyle=\color{string},
  backgroundcolor=\color{codebg},
  frame=single,
  rulecolor=\color{codeframe},
  frameround=tttt,
  breaklines=true,
  showstringspaces=false,
  tabsize=2,
  literate=
    {é}{{\'e}}1 {è}{{\`e}}1 {ê}{{\^e}}1 {à}{{\`a}}1 {â}{{\^a}}1
    {ç}{{\c{c}}}1 {²}{{$^2$}}1 {°}{{${}^\circ$}}1
    {α}{{$\alpha$}}1 {β}{{$\beta$}}1 {γ}{{$\gamma$}}1
    {δ}{{$\delta$}}1 {θ}{{$\theta$}}1 {λ}{{$\lambda$}}1
    {μ}{{$\mu$}}1 {σ}{{$\sigma$}}1 {ε}{{$\varepsilon$}}1
    {ρ}{{$\rho$}}1 {τ}{{$\tau$}}1 {φ}{{$\phi$}}1
    {≤}{{$\leq$}}1 {≥}{{$\geq$}}1 {≈}{{$\approx$}}1
}
\lstdefinestyle{Rstyle}{
  language=R,
  basicstyle=\ttfamily\scriptsize,
  keywordstyle=\color{keyword}\bfseries,
  commentstyle=\color{comment}\itshape,
  stringstyle=\color{string},
  backgroundcolor=\color{codebg},
  frame=single,
  rulecolor=\color{codeframe},
  frameround=tttt,
  breaklines=true,
  showstringspaces=false,
  tabsize=2
}
\lstset{style=pythonstyle}

% --- Common macros ---
% =========================================================
% Common math macros — Time Series Analysis module
% Author : Aymen Ben Brik (aymen.benbrik@esprit.tn)
% Esprit School of Business
% =========================================================

% --- Standard sets ---
\providecommand{\R}{\mathbb{R}}
\providecommand{\N}{\mathbb{N}}
\providecommand{\Z}{\mathbb{Z}}
\providecommand{\Q}{\mathbb{Q}}
\providecommand{\C}{\mathbb{C}}

% --- Probability / statistics ---
\providecommand{\E}{\mathbb{E}}
\providecommand{\Prob}{\mathbb{P}}
\providecommand{\Var}{\mathrm{Var}}
\providecommand{\Cov}{\mathrm{Cov}}
\providecommand{\Corr}{\mathrm{Corr}}
\providecommand{\indep}{\perp\!\!\!\perp}
\providecommand{\iid}{\overset{\text{i.i.d.}}{\sim}}
\providecommand{\given}{\,\big\vert\,}

% --- Estimators / hat ---
\providecommand{\hatm}{\widehat{m}}
\providecommand{\hatsig}{\widehat{\sigma}}
\providecommand{\hatrho}{\widehat{\rho}}
\providecommand{\hatphi}{\widehat{\phi}}
\providecommand{\hattheta}{\widehat{\theta}}

% --- Time series notation ---
\providecommand{\Lop}{L}                      % lag operator
\providecommand{\Diff}{\Delta}                % difference operator
\providecommand{\AR}[1]{\mathrm{AR}(#1)}
\providecommand{\MAm}[1]{\mathrm{MA}(#1)}
\providecommand{\ARMA}[2]{\mathrm{ARMA}(#1,#2)}
\providecommand{\ARIMA}[3]{\mathrm{ARIMA}(#1,#2,#3)}
\providecommand{\SARIMA}[6]{\mathrm{SARIMA}(#1,#2,#3)(#4,#5,#6)}
\providecommand{\ARCH}[1]{\mathrm{ARCH}(#1)}
\providecommand{\GARCH}[2]{\mathrm{GARCH}(#1,#2)}
\providecommand{\WN}{\mathrm{WN}}              % white noise
\providecommand{\MSE}{\mathrm{MSE}}
\providecommand{\MAE}{\mathrm{MAE}}
\providecommand{\RMSE}{\mathrm{RMSE}}
\providecommand{\AIC}{\mathrm{AIC}}
\providecommand{\BIC}{\mathrm{BIC}}
\providecommand{\ACF}{\mathrm{ACF}}
\providecommand{\PACF}{\mathrm{PACF}}

% --- Linear algebra ---
\providecommand{\transp}[1]{#1^{\!\top}}
\providecommand{\norm}[1]{\left\lVert #1 \right\rVert}
\providecommand{\abs}[1]{\left\lvert #1 \right\rvert}
\providecommand{\inner}[2]{\langle #1, #2 \rangle}
\providecommand{\diag}{\mathrm{diag}}
\providecommand{\tr}{\mathrm{tr}}

% --- Bold vectors / matrices ---
\providecommand{\vx}{\bm{x}}
\providecommand{\vy}{\bm{y}}
\providecommand{\vz}{\bm{z}}
\providecommand{\vh}{\bm{h}}
\providecommand{\vu}{\bm{u}}
\providecommand{\vv}{\bm{v}}
\providecommand{\vw}{\bm{w}}
\providecommand{\vb}{\bm{b}}
\providecommand{\vW}{\bm{W}}
\providecommand{\vSigma}{\bm{\Sigma}}
\providecommand{\veps}{\bm{\varepsilon}}

% --- Author identity (reused on titlepages) ---
\providecommand{\auteurNom}{Aymen Ben Brik}
\providecommand{\auteurEmail}{aymen.benbrik@esprit.tn}
\providecommand{\auteurInstitution}{Esprit School of Business}
\providecommand{\auteurRole}{Module head}
\providecommand{\moduleNom}{Time Series Analysis}
\providecommand{\anneeUniversitaire}{2025--2026}


% --- tcolorbox boxes for slides ---
\definecolor{boxTh}{RGB}{120, 30, 30}
\definecolor{boxProp}{RGB}{30, 100, 60}
\definecolor{boxDef}{RGB}{30, 60, 110}
\definecolor{boxEx}{RGB}{180, 120, 0}

\newtcolorbox{infobox}[1][]{enhanced, colback=espritBlue!5, colframe=espritBlue,
  fonttitle=\bfseries\small,
  title={\ifstrempty{#1}{}{#1}},
  boxrule=0.4pt, arc=2pt, left=4pt, right=4pt, top=2pt, bottom=2pt}
\newtcolorbox{warnbox}[1][]{enhanced, colback=espritAccent!5, colframe=espritAccent,
  fonttitle=\bfseries\small,
  title={\ifstrempty{#1}{}{#1}},
  boxrule=0.4pt, arc=2pt, left=4pt, right=4pt, top=2pt, bottom=2pt}
\newtcolorbox{thbox}[1][]{enhanced, colback=boxTh!5, colframe=boxTh,
  fonttitle=\bfseries\small,
  title={Theorem\ifstrempty{#1}{}{ -- #1}},
  boxrule=0.5pt, arc=2pt, left=4pt, right=4pt, top=2pt, bottom=2pt}
\newtcolorbox{propbox}[1][]{enhanced, colback=boxProp!5, colframe=boxProp,
  fonttitle=\bfseries\small,
  title={Proposition\ifstrempty{#1}{}{ -- #1}},
  boxrule=0.5pt, arc=2pt, left=4pt, right=4pt, top=2pt, bottom=2pt}
\newtcolorbox{defbox}[1][]{enhanced, colback=boxDef!5, colframe=boxDef,
  fonttitle=\bfseries\small,
  title={Definition\ifstrempty{#1}{}{ -- #1}},
  boxrule=0.5pt, arc=2pt, left=4pt, right=4pt, top=2pt, bottom=2pt}
\newtcolorbox{exbox}[1][]{enhanced, colback=boxEx!5, colframe=boxEx,
  fonttitle=\bfseries\small,
  title={Example\ifstrempty{#1}{}{ -- #1}},
  boxrule=0.5pt, arc=2pt, left=4pt, right=4pt, top=2pt, bottom=2pt}

% AFTER the line \documentclass[aspectratio=169]{beamer}.
% =========================================================

\usepackage[utf8]{inputenc}
\usepackage[T1]{fontenc}
\usepackage[english]{babel}
\usepackage{lmodern}
\usepackage{amsmath, amssymb, mathtools, bm}
\usepackage{graphicx}
\usepackage{booktabs}
\usepackage{array}
\usepackage{multirow}
\usepackage{colortbl}
\usepackage{xcolor}
\usepackage{tikz}
\usetikzlibrary{positioning, arrows.meta, calc, shapes, fit, backgrounds, matrix, decorations.pathreplacing}
\usepackage{listings}
\usepackage[most]{tcolorbox}

% --- Theme ---
\usetheme{metropolis}
\setbeamertemplate{navigation symbols}{}
\metroset{block=fill}

% --- Esprit colour palette ---
\definecolor{espritBlue}{RGB}{30, 60, 110}
\definecolor{espritAccent}{RGB}{200, 80, 30}
\setbeamercolor{progress bar}{fg=espritAccent}
\setbeamercolor{title separator}{fg=espritAccent}
\setbeamercolor{frametitle}{bg=espritBlue}

% --- Code listings (Python + R, with UTF-8 literate) ---
\definecolor{codebg}{RGB}{248,248,248}
\definecolor{codeframe}{RGB}{220,220,220}
\definecolor{keyword}{RGB}{0,82,155}
\definecolor{comment}{RGB}{88,110,117}
\definecolor{string}{RGB}{133,153,0}

\lstdefinestyle{pythonstyle}{
  language=Python,
  basicstyle=\ttfamily\scriptsize,
  keywordstyle=\color{keyword}\bfseries,
  commentstyle=\color{comment}\itshape,
  stringstyle=\color{string},
  backgroundcolor=\color{codebg},
  frame=single,
  rulecolor=\color{codeframe},
  frameround=tttt,
  breaklines=true,
  showstringspaces=false,
  tabsize=2,
  literate=
    {é}{{\'e}}1 {è}{{\`e}}1 {ê}{{\^e}}1 {à}{{\`a}}1 {â}{{\^a}}1
    {ç}{{\c{c}}}1 {²}{{$^2$}}1 {°}{{${}^\circ$}}1
    {α}{{$\alpha$}}1 {β}{{$\beta$}}1 {γ}{{$\gamma$}}1
    {δ}{{$\delta$}}1 {θ}{{$\theta$}}1 {λ}{{$\lambda$}}1
    {μ}{{$\mu$}}1 {σ}{{$\sigma$}}1 {ε}{{$\varepsilon$}}1
    {ρ}{{$\rho$}}1 {τ}{{$\tau$}}1 {φ}{{$\phi$}}1
    {≤}{{$\leq$}}1 {≥}{{$\geq$}}1 {≈}{{$\approx$}}1
}
\lstdefinestyle{Rstyle}{
  language=R,
  basicstyle=\ttfamily\scriptsize,
  keywordstyle=\color{keyword}\bfseries,
  commentstyle=\color{comment}\itshape,
  stringstyle=\color{string},
  backgroundcolor=\color{codebg},
  frame=single,
  rulecolor=\color{codeframe},
  frameround=tttt,
  breaklines=true,
  showstringspaces=false,
  tabsize=2
}
\lstset{style=pythonstyle}

% --- Common macros ---
% =========================================================
% Common math macros — Time Series Analysis module
% Author : Aymen Ben Brik (aymen.benbrik@esprit.tn)
% Esprit School of Business
% =========================================================

% --- Standard sets ---
\providecommand{\R}{\mathbb{R}}
\providecommand{\N}{\mathbb{N}}
\providecommand{\Z}{\mathbb{Z}}
\providecommand{\Q}{\mathbb{Q}}
\providecommand{\C}{\mathbb{C}}

% --- Probability / statistics ---
\providecommand{\E}{\mathbb{E}}
\providecommand{\Prob}{\mathbb{P}}
\providecommand{\Var}{\mathrm{Var}}
\providecommand{\Cov}{\mathrm{Cov}}
\providecommand{\Corr}{\mathrm{Corr}}
\providecommand{\indep}{\perp\!\!\!\perp}
\providecommand{\iid}{\overset{\text{i.i.d.}}{\sim}}
\providecommand{\given}{\,\big\vert\,}

% --- Estimators / hat ---
\providecommand{\hatm}{\widehat{m}}
\providecommand{\hatsig}{\widehat{\sigma}}
\providecommand{\hatrho}{\widehat{\rho}}
\providecommand{\hatphi}{\widehat{\phi}}
\providecommand{\hattheta}{\widehat{\theta}}

% --- Time series notation ---
\providecommand{\Lop}{L}                      % lag operator
\providecommand{\Diff}{\Delta}                % difference operator
\providecommand{\AR}[1]{\mathrm{AR}(#1)}
\providecommand{\MAm}[1]{\mathrm{MA}(#1)}
\providecommand{\ARMA}[2]{\mathrm{ARMA}(#1,#2)}
\providecommand{\ARIMA}[3]{\mathrm{ARIMA}(#1,#2,#3)}
\providecommand{\SARIMA}[6]{\mathrm{SARIMA}(#1,#2,#3)(#4,#5,#6)}
\providecommand{\ARCH}[1]{\mathrm{ARCH}(#1)}
\providecommand{\GARCH}[2]{\mathrm{GARCH}(#1,#2)}
\providecommand{\WN}{\mathrm{WN}}              % white noise
\providecommand{\MSE}{\mathrm{MSE}}
\providecommand{\MAE}{\mathrm{MAE}}
\providecommand{\RMSE}{\mathrm{RMSE}}
\providecommand{\AIC}{\mathrm{AIC}}
\providecommand{\BIC}{\mathrm{BIC}}
\providecommand{\ACF}{\mathrm{ACF}}
\providecommand{\PACF}{\mathrm{PACF}}

% --- Linear algebra ---
\providecommand{\transp}[1]{#1^{\!\top}}
\providecommand{\norm}[1]{\left\lVert #1 \right\rVert}
\providecommand{\abs}[1]{\left\lvert #1 \right\rvert}
\providecommand{\inner}[2]{\langle #1, #2 \rangle}
\providecommand{\diag}{\mathrm{diag}}
\providecommand{\tr}{\mathrm{tr}}

% --- Bold vectors / matrices ---
\providecommand{\vx}{\bm{x}}
\providecommand{\vy}{\bm{y}}
\providecommand{\vz}{\bm{z}}
\providecommand{\vh}{\bm{h}}
\providecommand{\vu}{\bm{u}}
\providecommand{\vv}{\bm{v}}
\providecommand{\vw}{\bm{w}}
\providecommand{\vb}{\bm{b}}
\providecommand{\vW}{\bm{W}}
\providecommand{\vSigma}{\bm{\Sigma}}
\providecommand{\veps}{\bm{\varepsilon}}

% --- Author identity (reused on titlepages) ---
\providecommand{\auteurNom}{Aymen Ben Brik}
\providecommand{\auteurEmail}{aymen.benbrik@esprit.tn}
\providecommand{\auteurInstitution}{Esprit School of Business}
\providecommand{\auteurRole}{Module head}
\providecommand{\moduleNom}{Time Series Analysis}
\providecommand{\anneeUniversitaire}{2025--2026}


% --- tcolorbox boxes for slides ---
\definecolor{boxTh}{RGB}{120, 30, 30}
\definecolor{boxProp}{RGB}{30, 100, 60}
\definecolor{boxDef}{RGB}{30, 60, 110}
\definecolor{boxEx}{RGB}{180, 120, 0}

\newtcolorbox{infobox}[1][]{enhanced, colback=espritBlue!5, colframe=espritBlue,
  fonttitle=\bfseries\small,
  title={\ifstrempty{#1}{}{#1}},
  boxrule=0.4pt, arc=2pt, left=4pt, right=4pt, top=2pt, bottom=2pt}
\newtcolorbox{warnbox}[1][]{enhanced, colback=espritAccent!5, colframe=espritAccent,
  fonttitle=\bfseries\small,
  title={\ifstrempty{#1}{}{#1}},
  boxrule=0.4pt, arc=2pt, left=4pt, right=4pt, top=2pt, bottom=2pt}
\newtcolorbox{thbox}[1][]{enhanced, colback=boxTh!5, colframe=boxTh,
  fonttitle=\bfseries\small,
  title={Theorem\ifstrempty{#1}{}{ -- #1}},
  boxrule=0.5pt, arc=2pt, left=4pt, right=4pt, top=2pt, bottom=2pt}
\newtcolorbox{propbox}[1][]{enhanced, colback=boxProp!5, colframe=boxProp,
  fonttitle=\bfseries\small,
  title={Proposition\ifstrempty{#1}{}{ -- #1}},
  boxrule=0.5pt, arc=2pt, left=4pt, right=4pt, top=2pt, bottom=2pt}
\newtcolorbox{defbox}[1][]{enhanced, colback=boxDef!5, colframe=boxDef,
  fonttitle=\bfseries\small,
  title={Definition\ifstrempty{#1}{}{ -- #1}},
  boxrule=0.5pt, arc=2pt, left=4pt, right=4pt, top=2pt, bottom=2pt}
\newtcolorbox{exbox}[1][]{enhanced, colback=boxEx!5, colframe=boxEx,
  fonttitle=\bfseries\small,
  title={Example\ifstrempty{#1}{}{ -- #1}},
  boxrule=0.5pt, arc=2pt, left=4pt, right=4pt, top=2pt, bottom=2pt}


\title[\moduleNom\ -- Ch.3]{Stationary Time Series}
\subtitle{Definitions $\cdot$ Autocovariance $\cdot$ ACF \& PACF $\cdot$ Wold}
\author{\auteurNom \\ \footnotesize\auteurEmail}
\institute{\auteurInstitution}
\date{\anneeUniversitaire}

\begin{document}

\begin{frame}\titlepage\end{frame}

\begin{frame}{Learning objectives}
  By the end of this chapter, you will be able to:
  \begin{itemize}
    \item state the definitions of \emph{strict} and \emph{weak (covariance)
          stationarity};
    \item compute and interpret the \emph{autocovariance}
          $\gamma(h)$ and \emph{autocorrelation} $\rho(h)$ functions;
    \item read the \emph{ACF} and \emph{PACF} of an empirical series
          and identify candidate AR/MA orders;
    \item state the \emph{Wold decomposition} and its consequences;
    \item run a \emph{white-noise test} (Box--Pierce, Ljung--Box) to
          assess residuals.
  \end{itemize}
\end{frame}

\begin{frame}{Chapter outline}
  \begin{enumerate}
    \item Strict vs weak stationarity
    \item Autocovariance and autocorrelation
    \item Partial autocorrelation (PACF)
    \item ACF / PACF signatures of standard processes
    \item Wold decomposition
    \item White-noise tests
  \end{enumerate}
\end{frame}

% =====================================================================
\section{Stationarity}

\begin{frame}{1.~Why stationarity?}
  Most classical time-series tools (estimation, forecasting,
  inference) require the data-generating process to have
  \emph{stable statistical properties} over time. Without
  stationarity:
  \begin{itemize}
    \item the sample mean and variance lose their interpretation;
    \item ACF estimates become inconsistent;
    \item forecasts diverge and confidence intervals are wrong.
  \end{itemize}
  \medskip
  \begin{infobox}
    The decomposition step (Chapter 2) extracts trend and seasonality;
    the residual is what we hope is stationary.
  \end{infobox}
\end{frame}

\begin{frame}{1.~Strict (strong) stationarity}
  \begin{defbox}[Strict stationarity]
    A process $\{X_t\}_{t \in \Z}$ is \emph{strictly stationary} if for
    every integer $k$, every $t_1 < \dots < t_k$ and every shift $h$,
    the joint distribution
    \[
      (X_{t_1}, X_{t_2}, \dots, X_{t_k})
      \overset{\text{d}}{=}
      (X_{t_1+h}, X_{t_2+h}, \dots, X_{t_k+h}).
    \]
  \end{defbox}
  \medskip
  \begin{itemize}
    \item Strong: \emph{all} finite-dimensional distributions are time-invariant.
    \item Implies stationarity of all moments (when they exist).
    \item Strict but rare; in practice we work with the weaker version.
  \end{itemize}
\end{frame}

\begin{frame}{1.~Weak (covariance) stationarity}
  \begin{defbox}[Weak stationarity]
    A process $\{X_t\}$ with $\E[X_t^2] < \infty$ is \emph{weakly
    stationary} if there exist $\mu \in \R$ and $\gamma : \Z \to \R$ such that
    \[
      \E[X_t] = \mu, \qquad \Cov(X_t, X_{t+h}) = \gamma(h)
      \quad \text{for all } t, h \in \Z.
    \]
  \end{defbox}
  \medskip
  \begin{itemize}
    \item Constant mean, constant variance ($\gamma(0)$), and
          autocovariance depending only on the lag $h$.
    \item Strict $\Rightarrow$ weak (under finite second moments);
          the converse fails in general (it holds for Gaussian processes).
  \end{itemize}
\end{frame}

\begin{frame}{1.~Strict vs weak stationarity}
  \begin{center}
  \begin{tabular}{lcc}
    \toprule
    \textbf{Property}                & \textbf{Strict} & \textbf{Weak} \\
    \midrule
    Joint distribution invariant     & yes             & not required \\
    Mean constant                    & yes             & yes          \\
    Variance constant                & yes             & yes          \\
    $\gamma(h)$ depends only on $h$  & yes             & yes          \\
    Higher moments invariant         & yes             & not required \\
    Useful for AR / MA / ARMA        & yes (overkill)  & yes          \\
    \bottomrule
  \end{tabular}
  \end{center}
  \medskip
  \begin{infobox}
    Throughout this course, ``stationary'' means \emph{weakly
    stationary} unless otherwise specified.
  \end{infobox}
\end{frame}

\begin{frame}{1.~White noise: the canonical stationary process}
  \begin{defbox}[White noise]
    $(\varepsilon_t)$ is a \emph{white noise} of variance $\sigma^2$
    (denoted $\WN(0, \sigma^2)$) if:
    \[
      \E[\varepsilon_t] = 0, \quad \Var[\varepsilon_t] = \sigma^2,
      \quad \Cov(\varepsilon_t, \varepsilon_{t+h}) = 0 \text{ for all } h \neq 0.
    \]
    Stronger: \emph{strong} white noise if the $\varepsilon_t$ are i.i.d.
  \end{defbox}
  \medskip
  \begin{itemize}
    \item ACF: $\rho(0) = 1$, $\rho(h) = 0$ for $h \neq 0$.
    \item No structure to exploit. Used as input to AR/MA models.
    \item In Python: \texttt{np.random.normal(0, sigma, n)} (Gaussian WN).
  \end{itemize}
\end{frame}

\begin{frame}{1.~Random walk: a canonical \emph{non}-stationary process}
  \begin{defbox}[Random walk]
    $X_t = X_{t-1} + \varepsilon_t$ with $\varepsilon_t \sim \WN(0, \sigma^2)$
    and $X_0 = 0$. Equivalently $X_t = \sum_{s=1}^{t}\varepsilon_s$.
  \end{defbox}
  \medskip
  \textbf{Why not stationary?}
  \[
    \Var[X_t] = t\, \sigma^2 \;\to\; \infty,
    \qquad
    \Cov(X_t, X_{t+h}) = t\, \sigma^2 \quad (\text{depends on } t).
  \]
  \medskip
  \begin{warnbox}
    Variance grows linearly with $t$. Many financial price series
    behave like a random walk; Chapter 4 introduces tools to test
    for it (ADF, KPSS).
  \end{warnbox}
\end{frame}

% =====================================================================
\section{Autocovariance and ACF}

\begin{frame}{2.~Autocovariance and autocorrelation}
  \begin{defbox}[ACVF and ACF]
    For a weakly stationary process with mean $\mu$:
    \[
      \gamma(h) \;=\; \Cov(X_t, X_{t+h}) = \E[(X_t - \mu)(X_{t+h} - \mu)],
    \]
    \[
      \rho(h) \;=\; \frac{\gamma(h)}{\gamma(0)} \;\in\; [-1, 1].
    \]
  \end{defbox}
  \medskip
  \textbf{Properties.}
  \begin{itemize}
    \item $\gamma(0) = \Var[X_t] \geq 0$.
    \item $\abs{\gamma(h)} \leq \gamma(0)$ (Cauchy--Schwarz).
    \item $\gamma(-h) = \gamma(h)$ — even function.
    \item $\gamma$ is a \emph{positive-definite} function.
  \end{itemize}
\end{frame}

\begin{frame}{2.~Sample ACF}
  Given observations $X_1, \dots, X_T$:
  \[
    \widehat{\gamma}(h) \;=\; \frac{1}{T}\sum_{t = 1}^{T - h}
      \bigl(X_t - \bar{X}\bigr)\bigl(X_{t+h} - \bar{X}\bigr),
    \qquad
    \widehat{\rho}(h) \;=\; \frac{\widehat{\gamma}(h)}{\widehat{\gamma}(0)}.
  \]
  \medskip
  \begin{thbox}[Bartlett's formula, white noise]
    Under $\WN(0, \sigma^2)$, for fixed $h \geq 1$:
    \[
      \sqrt{T}\,\widehat{\rho}(h) \;\xrightarrow{d}\; \mathcal{N}(0, 1).
    \]
  \end{thbox}
  \medskip
  \textbf{Practical bands.} On the ACF plot, draw $\pm 1.96/\sqrt{T}$
  bands; spikes outside are ``significant'' at the 5\% level.
\end{frame}

% =====================================================================
\section{Partial autocorrelation (PACF)}

\begin{frame}{3.~Partial autocorrelation}
  \begin{defbox}[PACF]
    $\alpha(h)$ is the correlation between $X_t$ and $X_{t+h}$
    \emph{after removing} the linear effect of $X_{t+1}, \dots, X_{t+h-1}$:
    \[
      \alpha(h) \;=\; \Corr\bigl(X_t - \widehat{X}_t,\; X_{t+h} - \widehat{X}_{t+h}\bigr),
    \]
    where $\widehat{X}_t$ and $\widehat{X}_{t+h}$ are the linear projections
    on $\{X_{t+1}, \dots, X_{t+h-1}\}$.
  \end{defbox}
  \medskip
  \begin{itemize}
    \item ACF: \emph{total} correlation (direct + indirect).
    \item PACF: \emph{direct} correlation only.
    \item Computed recursively (\emph{Durbin--Levinson} algorithm).
  \end{itemize}
\end{frame}

% =====================================================================
\section{ACF / PACF signatures}

\begin{frame}{4.~Theoretical ACF and PACF}
  \begin{center}
  \begin{tabular}{lll}
    \toprule
    \textbf{Process} & \textbf{ACF $\rho(h)$}             & \textbf{PACF $\alpha(h)$} \\
    \midrule
    White noise      & $\rho(h) = 0$ for $h \neq 0$       & $\alpha(h) = 0$ for $h \geq 1$ \\
    AR$(p)$          & decays exponentially / sinusoidally & cuts off at $h = p$       \\
    MA$(q)$          & cuts off at $h = q$                & decays exponentially       \\
    ARMA$(p, q)$     & decays after $h > q$               & decays after $h > p$       \\
    \bottomrule
  \end{tabular}
  \end{center}
  \medskip
  \begin{infobox}
    \emph{Box--Jenkins identification.} Plot the sample ACF and PACF;
    use the cutoff/decay patterns above to choose tentative $(p, q)$.
  \end{infobox}
\end{frame}

\begin{frame}{4.~AR(1) example: $X_t = \phi X_{t-1} + \varepsilon_t$, $\abs\phi < 1$}
  Theoretical autocorrelation:
  \[
    \rho(h) \;=\; \phi^{\,\abs{h}}.
  \]
  \medskip
  \begin{itemize}
    \item $\phi = 0.8$: $\rho(1) = 0.8$, $\rho(2) = 0.64$, $\rho(5) = 0.328$, $\rho(10) = 0.107$.
    \item $\phi = -0.8$: alternating sign, same magnitudes.
    \item PACF cuts off at lag 1: $\alpha(1) = \phi$, $\alpha(h) = 0$
          for $h \geq 2$.
  \end{itemize}
  \medskip
  \begin{warnbox}
    Stationarity of AR(1) requires $\abs\phi < 1$. If $\phi = 1$, we
    obtain the random walk (non-stationary).
  \end{warnbox}
\end{frame}

\begin{frame}{4.~MA(1) example: $X_t = \varepsilon_t + \theta \varepsilon_{t-1}$}
  Theoretical autocorrelation:
  \[
    \rho(0) = 1,
    \quad
    \rho(1) = \frac{\theta}{1 + \theta^2},
    \quad
    \rho(h) = 0 \text{ for } h \geq 2.
  \]
  \medskip
  \begin{itemize}
    \item ACF \emph{cuts off} at lag 1: clean signature.
    \item PACF \emph{decays exponentially} (no cutoff).
    \item For $\theta = 0.5$: $\rho(1) = 0.4$.
  \end{itemize}
\end{frame}

% =====================================================================
\section{Wold decomposition}

\begin{frame}{5.~Wold decomposition (1938)}
  \begin{thbox}
    Every \emph{purely non-deterministic} weakly stationary process
    $X_t$ admits a unique decomposition
    \[
      X_t \;=\; \mu \;+\; \sum_{j = 0}^{\infty} \psi_j\, \varepsilon_{t - j},
      \qquad \sum_{j = 0}^{\infty} \psi_j^2 < \infty,
    \]
    with $\psi_0 = 1$ and $(\varepsilon_t)$ a white noise of variance
    $\sigma^2$ (the \emph{innovations}).
  \end{thbox}
  \medskip
  \begin{itemize}
    \item Existence and uniqueness: any stationary process is an
          \emph{infinite moving average} of innovations.
    \item Foundation of ARMA modelling: ARMA processes are
          finite-dimensional approximations of the Wold representation.
  \end{itemize}
\end{frame}

% =====================================================================
\section{White-noise tests}

\begin{frame}{6.~Box--Pierce and Ljung--Box tests}
  Test $H_0$: the residuals are white noise (no autocorrelation up to
  lag $m$).
  \medskip
  \begin{defbox}[Test statistics]
    \[
      Q_{\text{BP}} = T \sum_{h = 1}^{m} \widehat\rho(h)^2,
      \qquad
      Q_{\text{LB}} = T(T + 2) \sum_{h = 1}^{m} \frac{\widehat\rho(h)^2}{T - h}.
    \]
    Under $H_0$, $Q_{\text{LB}} \overset{a}{\sim} \chi^2_{m - p - q}$
    where $p + q$ is the number of estimated ARMA parameters
    (or simply $m$ if no model has been fitted).
  \end{defbox}
  \medskip
  \begin{itemize}
    \item Reject $H_0$ if $Q_{\text{LB}} > \chi^2_{m - p - q,\, 1 - \alpha}$.
    \item Choose $m$ a few (10--40); for monthly data, $m \in \{12, 24, 36\}$.
  \end{itemize}
\end{frame}

\begin{frame}[fragile]{6.~Ljung--Box in Python}
\begin{lstlisting}
from statsmodels.stats.diagnostic import acorr_ljungbox

# residuals after decomposition / ARMA fitting
out = acorr_ljungbox(residuals, lags=[12, 24, 36], return_df=True)
print(out)
#       lb_stat  lb_pvalue
# 12     ...     ...
# 24     ...     ...
# 36     ...     ...
\end{lstlisting}
  \medskip
  Small $p$-value ($< 0.05$) $\Rightarrow$ residuals are \emph{not}
  white noise: there is structure left to model.
\end{frame}

\begin{frame}{Up next}
  \textbf{Chapter 4.} Non-stationary time series: random walk,
  unit-root tests (ADF, KPSS), differencing, ARIMA.
  \medskip
  \begin{itemize}
    \item How to formally test for stationarity.
    \item Transformations to achieve stationarity (differencing, log).
    \item Integrated processes.
  \end{itemize}
  \medskip
  \begin{center}\Large\bfseries Thank you.\end{center}
  \begin{center}\footnotesize\auteurNom\ -- \auteurEmail\ -- \auteurInstitution\end{center}
\end{frame}

\end{document}
